%% =========================================================================
%% MODELO CANÔNICO ABNTEX2 PARA TRABALHOS ACADÊMICOS & RELATÓRIOS TÉCNICOS
%% Ecossistema Aluno Nota 90 — Ciência da Computação / Engenharia de Software
%% =========================================================================
\documentclass[
    12pt,               % Tamanho da fonte padrão ABNT
    openright,          % Capítulos começam em página ímpar
    oneside,            % Para impressão em apenas um lado da folha (ou twoside)
    a4paper,            % Tamanho do papel A4
    english,            % Idioma adicional para abstract
    brazil              % Idioma principal do documento
]{abntex2}

% --- Pacotes Fundamentais ---
\usepackage[utf8]{inputenc}        % Codificação do documento (letras acentuadas)
\usepackage[T1]{fontenc}           % Seleção de códigos de fonte
\usepackage{lmodern}              % Usa a fonte Latin Modern
\usepackage{color}                % Controle das cores
\usepackage{graphicx}             % Inclusão de gráficos e figuras
\usepackage{microtype}            % Melhorias na justificação do texto
\usepackage{amsmath,amssymb}      % Fórmulas matemáticas avançadas
\usepackage{booktabs}             % Tabelas de alta qualidade acadêmica
\usepackage{listings}             % Inclusão de código-fonte
\usepackage{xcolor}               % Suporte avançado a cores

% --- Pacotes de Citações ABNT (Autor-Data) ---
\usepackage[brazilian,hyperpageref]{backref}  % Páginas com as citações na bibliografia
\usepackage[alf,abnt-emphasize=bf,abnt-repeated-title-omission=yes]{abntex2cite} % Citações padrão ABNT

% --- Configurações do Pacote Listings (Código-Fonte com suporte a UTF-8) ---
\definecolor{codegreen}{rgb}{0,0.6,0}
\definecolor{codegray}{rgb}{0.5,0.5,0.5}
\definecolor{codepurple}{rgb}{0.58,0,0.82}
\definecolor{backcolour}{rgb}{0.96,0.96,0.96}

\lstdefinestyle{mystyle}{
    backgroundcolor=\color{backcolour},   
    commentstyle=\color{codegreen},
    keywordstyle=\color{blue}\bfseries,
    numberstyle=\tiny\color{codegray},
    stringstyle=\color{codepurple},
    basicstyle=\ttfamily\footnotesize,
    breakatwhitespace=false,         
    breaklines=true,                 
    captionpos=b,                    
    keepspaces=true,                 
    numbers=left,                    
    numbersep=5pt,                  
    showspaces=false,                
    showstringspaces=false,
    showtabs=false,                  
    tabsize=2
}
\lstset{style=mystyle}

% --- Informações do Documento ---
\titulo{Título do Trabalho Acadêmico: Subtítulo Explicativo}
\autor{Eduardo Faria}
\local{Brasil}
\data{2026}
\orientador{Prof. Dr. Nome do Orientador}
\instituicao{%
  Universidade de Computação e Tecnologia
  \par
  Curso de Engenharia de Software / Ciência da Computação}
\tipotrabalho{Trabalho Acadêmico / Relatório Técnico}
\preambulo{Trabalho apresentado como requisito parcial para obtenção de nota na disciplina X.}

% --- Configurações de hiperlinks ---
\definecolor{blue_abnt}{RGB}{41,5,195}
\hypersetup{
    colorlinks=true,
    linkcolor=blue_abnt,
    citecolor=blue_abnt,
    urlcolor=blue_abnt
}

% =========================================================================
% INÍCIO DO DOCUMENTO
% =========================================================================
\begin{document}
\selectlanguage{brazil}
\frenchspacing 

% --- Elementos Pré-Textuais ---
\imprimircapa
\imprimirfolhaderosto

% Resumo em Português
\setlength{\absparsep}{18pt}
\begin{resumo}
 Este documento apresenta o modelo estrutural canônico em conformidade com as normas da ABNT utilizando o pacote \textit{abntex2}. O objetivo deste trabalho é fornecer uma base sólida, reprodutível e livre de erros para a elaboração de artigos, relatórios técnicos e monografias nas disciplinas de computação e engenharia.
 
 \textbf{Palavras-chave}: Engenharia de Software. LaTeX. ABNT. Computação.
\end{resumo}

% Sumário
\pdfbookmark[0]{\contentsname}{toc}
\tableofcontents*
\cleardoublepage

% --- Elementos Textuais ---
\textual

\chapter{Introdução}
A introdução contextualiza o problema investigado, a motivação do estudo, a relevância técnica para a área de computação e os objetivos gerais e específicos da pesquisa.

Segundo \citeonline{knuth1984literate}, a programação e a documentação técnica devem andar de mãos dadas, promovendo clareza cognitiva e rigor estrutural.

\chapter{Fundamentação Teórica \& Metodologia}
Aqui são detalhados os conceitos teóricos, padrões arquiteturais e algoritmos envolvidos.

\section{Modelagem Matemática}
Equações complexas mantêm formatação impecável:
\begin{equation}
    f(x) = \sum_{i=1}^{n} \frac{\alpha_i \cdot \beta_i}{\sqrt{1 + \gamma^2}} + \int_{0}^{\infty} e^{-\lambda t} dt
\end{equation}

\section{Exemplo de Código-Fonte}
Exemplo de algoritmo implementado com realce de sintaxe:

\begin{lstlisting}[language=Python, caption={Exemplo de algoritmo de busca em Python.}]
def busca_binaria(vetor, elemento):
    """Executa busca binaria em O(log n)."""
    inicio = 0
    fim = len(vetor) - 1
    
    while inicio <= fim:
        meio = (inicio + fim) // 2
        if vetor[meio] == elemento:
            return meio
        elif vetor[meio] < elemento:
            inicio = meio + 1
        else:
            fim = meio - 1
    return -1
\end{lstlisting}

\chapter{Resultados e Discussão}
Apresentação dos dados empíricos, testes de benchmark, métricas de complexidade ciclomática e tempo de execução.

\chapter{Conclusão}
Síntese das contribuições alcançadas e proposição de trabalhos futuros.

% --- Elementos Pós-Textuais ---
\postextual

% Bibliografia
\bibliography{referencias}

\end{document}
